\chapter{Transformer}
\section{Seq2Seq模型}
语言模型根据已经出现的单词输出下一个出现的单词的概率分布，通过对这个概率分布进行采样得到下一个词，
再将已生成的词序列不断送回到模型中，重复该过程，直到生成结束符号 <eos> 为止，就可以生成完整的文本。

在上述过程中，如何选择下一个词取决于具体的采样策略。
最简单的策略是贪心搜索（Greedy Search），即每次选择概率最高的词作为下一个输出，
但这种方法容易生成重复或单调的文本。为了增加生成的多样性，可以采用以下策略：
\begin{itemize}
	\item \textbf{温度采样（Temperature Sampling）}：引入温度参数 $T$，将 softmax 的输入除以 $T$。
	      当 $T < 1$ 时，概率分布更尖锐，模型更倾向于选择高概率的词；
	      当 $T > 1$ 时，分布更平坦，生成的随机性更大。
	\item \textbf{Top-k 采样}：只保留概率最高的 $k$ 个词，将其余词的概率置为零后再进行采样。
	\item \textbf{Top-p 采样（Nucleus Sampling）}：动态选择一个最小的词集，使其累积概率之和达到阈值 $p$，
	      然后在该词集中进行采样。与 Top-k 不同，Top-p 的词集大小是动态变化的。
\end{itemize}

Seq2Seq模型则是另一类文本生成模型，用于将一种序列转换为另一种序列，广泛应用于机器翻译、语音识别等任务。
它通常由两个模块组成：Encoder和Decoder，因此也被称为 Encoder-Decoder 模型。
编码器常见结构是 Embedding 层加 RNN 层，用于将输入序列编码为一个向量表示；
解码器则根据这个向量表示逐步生成输出序列，其生成过程可以参考前面的采样策略。

\begin{figure}[ht]
	\centering
	\includegraphics[width=0.4\textwidth, keepaspectratio]{6-1.pdf}
	\caption{Seq2Seq模型示例}
\end{figure}

编码器的输出通常是一个固定长度的向量，但输入序列的长度是可变的，将所有信息压缩到一个固定向量中可能导致部分信息丢失。
一个自然的想法是保留编码器在每个时间步的隐藏状态，将它们拼接成一个矩阵，让解码器能够利用全部时间步的信息。
然而，解码器的输入是单个向量，因此需要先对矩阵的行向量进行加权平均，将其“压缩”为单个向量。
而如何确定加权平均中的权重，就是注意力机制（Attention Mechanism）所要解决的问题。

\section{Self-Attention}
在基于 RNN 的 Seq2Seq 模型中，注意力机制使得解码器在生成每个输出单词时，能够动态地关注输入序列中的不同部分。
例如，在将 “I love you” 翻译为 “我爱你” 的过程中，生成“我”时模型更倾向于关注“I”，而生成“爱”时则更多关注“love”。
既然输出序列可以使用注意力机制，那么输入序列是否也可以使用注意力机制呢？

答案是肯定的，这正是 Transformer 模型中自注意力机制（Self-Attention）的作用。
需要说明的是，Transformer 作为 Seq2Seq 模型的一种实现方式，也包含类似于基于 RNN 的 Seq2Seq 模型中的注意力机制，
这部分在 Transformer 中被称为交叉注意力（Cross-Attention）。
让我们通过一个例子来快速理解自注意力机制，例句：
\begin{quotation}
	A dog ate the food because it was hungry
\end{quotation}
显然，例句中的代词 it 指代的是 dog  而不是 food，自注意力机制的作用就是关注句子间单词的关系。

将句子输入模型的编码器，得到一个输出矩阵 $\mathbf{X}$，为了计算自注意力，我们需要再创建三个新的矩阵：
查询（query）矩阵 $\mathbf{Q}$、键（key）矩阵 $\mathbf{K}$，以及值（value）矩阵 $\mathbf{V}$。为了创建它们，
我们需要先创建另外三个权重（weight）矩阵，$\mathbf{W}^Q$、$\mathbf{W}^K$、$\mathbf{W}^V$。
用矩阵 $\mathbf{X}$ 分别乘以矩阵 $\mathbf{W}^Q$、$\mathbf{W}^K$、$\mathbf{W}^V$，
依次创建出 $\mathbf{Q}$、 $\mathbf{K}$ 和 $\mathbf{V}$。
值得注意的是，权重矩阵 $\mathbf{W}^Q$、$\mathbf{W}^K$ 和 $\mathbf{W}^V$ 的初始值完全是随机的，
但最优值则需要通过训练获得。我们取得的权值越优，通过计算所得的查询矩阵、键矩阵和值矩阵也会越精确。

假设输入句子为“I am good”，自注意力机制首先要计算查询矩阵 $\boldsymbol{Q}$ 与键矩阵 $\boldsymbol{K}^\mathrm{T}$ 的点积。
$$
	\boldsymbol{Q} \cdot \boldsymbol{K}^\mathrm{T} =
	\begin{bmatrix}
		\boldsymbol{q}_1 \cdot \boldsymbol{k}_1 & \boldsymbol{q}_1 \cdot \boldsymbol{k}_2 & \boldsymbol{q}_1 \cdot \boldsymbol{k}_3 \\
		\boldsymbol{q}_2 \cdot \boldsymbol{k}_1 & \boldsymbol{q}_2 \cdot \boldsymbol{k}_2 & \boldsymbol{q}_2 \cdot \boldsymbol{k}_3 \\
		\boldsymbol{q}_3 \cdot \boldsymbol{k}_1 & \boldsymbol{q}_3 \cdot \boldsymbol{k}_2 & \boldsymbol{q}_3 \cdot \boldsymbol{k}_3 \\
	\end{bmatrix}
$$

首先，来看 $\boldsymbol{Q} \cdot \boldsymbol{K}^\mathrm{T}$ 矩阵的第一行，可以看到，
这一行计算的是查询向量 $\boldsymbol{q}_1$（I）与所有的键向量 $\boldsymbol{k}_1$（I）、$\boldsymbol{k}_2$（am）和 $\boldsymbol{k}_3$（good）的点积。
通过计算两个向量的点积可以知道它们之间的相似度。
因此，通过计算查询向量（$\boldsymbol{q}_1$）和键向量（$\boldsymbol{k}_1$、$\boldsymbol{k}_2$、$\boldsymbol{k}_3$）的点积，
可以了解单词 I 与句子中的所有单词的相似度。我们了解到, I 这个词与自己的关系比与 am 和 good 这两个词的关系更紧密，
因为点积值 110 大于 90 和 80。

\begin{table}[ht]
	\centering
	$$
		\boldsymbol{Q} \cdot \boldsymbol{K}^\mathrm{T} =
		\begin{bmatrix}
			110 & 90 & 80  \\
			70  & 99 & 70  \\
			90  & 70 & 100 \\
		\end{bmatrix}
	$$
	\caption{点积示例，本节的结果均为假设值，后续不再说明}
\end{table}

自注意力机制的第 2 步是将 $\boldsymbol{Q} \cdot \boldsymbol{K}^\mathrm{T}$ 矩阵除以键向量维度的平方根 $\sqrt{d_k}$。
这是因为当 $d_k$ 较大时，点积结果的方差也会较大，使得 softmax 的输入值差异悬殊，输出接近 one-hot 分布，导致梯度消失。
除以 $\sqrt{d_k}$ 可以将点积的方差重新缩放回去，从而获得更稳定的梯度。

目前所得的相似度分数尚未被归一化，所以第 3 步是使用 softmax 函数对其进行归一化处理。
应用 softmax 函数将使数值分布在 0 到 1 的范围内，且每一行的所有数之和等于 1。
这个矩阵被称为分数（score）矩阵，通过它，我们可以了解句子中的每个词与所有词的相关程度。
比如：I 这个词与它本身的相关程度是 90\%，与 am 这个词的相关程度是 7\%，与 good 这个词的相关程度是 3\%。
$$
	\text{softmax}\left(\frac{\boldsymbol{Q} \cdot \boldsymbol{K}^\mathrm{T}}{\sqrt{d_k}}\right) =
	\begin{bmatrix}
		            & \text{I} & \text{am} & \text{good} \\
		\text{I}    & 0.90     & 0.07      & 0.03        \\
		\text{am}   & 0.025    & 0.95      & 0.025       \\
		\text{good} & 0.21     & 0.03      & 0.76        \\
	\end{bmatrix}
$$

自注意力机制的最后一步是计算注意力矩阵 $\boldsymbol{Z}$，
注意力矩阵包含句子中每个单词的注意力向量。它通过将分数矩阵乘以值矩阵 $\boldsymbol{V}$ 得出。

\begin{align*}
	Z & = \text{softmax}\left(\frac{\boldsymbol{Q} \cdot \boldsymbol{K}^\mathrm{T}}{\sqrt{d_k}}\right) \boldsymbol{V} \\
	  & =
	\begin{bmatrix}
		0.90  & 0.07 & 0.03  \\
		0.025 & 0.95 & 0.025 \\
		0.21  & 0.03 & 0.76  \\
	\end{bmatrix}
	\begin{bmatrix}
		67.85 & 91.2 & \ldots & 0.13 \\
		13.13 & 63.1 & \ldots & 4.44 \\
		12.12 & 96.1 & \ldots & 43.4 \\
	\end{bmatrix}                                                                                      \\                                                                                              \\
	  & =
	\begin{bmatrix}
		\boldsymbol{z}_1 \\
		\boldsymbol{z}_2 \\
		\boldsymbol{z}_3 \\
	\end{bmatrix}
\end{align*}

可以看出，单词 I 的自注意力向量 $\boldsymbol{z}_1$ 是分数加权的值向量之和。
所以，$\boldsymbol{z}_1$ 的值将包含 90\% 的值向量 $\boldsymbol{v}_1$（I）、7\% 的值向量 $\boldsymbol{v}_2$（am），
以及 3\% 的值向量 $\boldsymbol{v}_3$（good）。
$$
	\boldsymbol{z}_1 = 0.90 \begin{bmatrix} 67.85 & 91.2 & \ldots \end{bmatrix} + 0.07 \begin{bmatrix} 13.13 & 63.1 & \ldots \end{bmatrix} + 0.03 \begin{bmatrix} 12.12 & 96.1 & \ldots \end{bmatrix}
$$
\begin{figure}[ht]
	\centering
	\includegraphics[width=0.2\textwidth, keepaspectratio]{6-2.pdf}
	\caption{自注意力计算流程}
\end{figure}

\section{位置编码}
在 RNN 中，句子是逐词送入网络的，以 I am good 为例，首先把 I 作为输入，接下来是 am，最后是good。
然而，Transformer 并不遵循这个模式，它将句子中的所有词并行地输入到神经网络中。
这种方式下，词序需要通过其他方式表示，Transformer 所采用的方式是位置编码。

在 Transformer 中，位置编码以矩阵的形式表示，位置编码矩阵 $\boldsymbol{P}$ 的维度与输入矩阵 $\boldsymbol{X}$ 的维度相同。
在将输入矩阵直接传给 Transformer 之前，我们需将位置编码矩阵 $\boldsymbol{P}$ 加到输入矩阵 $\boldsymbol{X}$ 中，再将其作为输入送入网络。
Transformer 中使用了正弦函数来计算位置编码：
\begin{equation*}
	P(\text{pos}, 2i) = \sin (\frac{\text{pos}}{10000^{2i/d_{\text{model}}}}) \quad
	P(\text{pos}, 2i+1) = \cos (\frac{\text{pos}}{10000^{2i/d_{\text{model}}}})
\end{equation*}

在上面的等式中，$\text{pos}$ 表示该词在句子中的位置，$i$ 表示在输入矩阵中的位置。
$$
	\boldsymbol{P} =
	\begin{bmatrix}
		\sin\left(\frac{\text{pos}}{10000^0}\right) & \cos\left(\frac{\text{pos}}{10000^0}\right) & \sin\left(\frac{\text{pos}}{10000^{2/4}}\right) & \cos\left(\frac{\text{pos}}{10000^{2/4}}\right) \\
		\sin\left(\frac{\text{pos}}{10000^0}\right) & \cos\left(\frac{\text{pos}}{10000^0}\right) & \sin\left(\frac{\text{pos}}{10000^{2/4}}\right) & \cos\left(\frac{\text{pos}}{10000^{2/4}}\right) \\
		\sin\left(\frac{\text{pos}}{10000^0}\right) & \cos\left(\frac{\text{pos}}{10000^0}\right) & \sin\left(\frac{\text{pos}}{10000^{2/4}}\right) & \cos\left(\frac{\text{pos}}{10000^{2/4}}\right) \\
	\end{bmatrix}
$$

上述方法属于绝对位置编码，它将位置信息作为一个固定的向量加到输入上。
另一种广泛使用的位置编码方法是旋转位置编码（Rotary Position Embedding，RoPE）。
RoPE 通过对查询向量和键向量施加一个与位置相关的旋转，使得两个向量的点积自然地编码了它们之间的相对位置信息。
具体来说，RoPE 将向量的相邻两个维度视为一个二维平面上的点，然后根据位置 $m$ 旋转一个角度 $m\theta$：
$$
	\begin{bmatrix} q_{2i}' \\ q_{2i+1}' \end{bmatrix}
	=
	\begin{bmatrix} \cos m\theta_i & -\sin m\theta_i \\ \sin m\theta_i & \cos m\theta_i \end{bmatrix}
	\begin{bmatrix} q_{2i} \\ q_{2i+1} \end{bmatrix}
$$
其中 $\theta_i = 10000^{-2i/d}$。
经过旋转后，$\boldsymbol{q}_m$ 与 $\boldsymbol{k}_n$ 的点积只依赖于相对位置 $m - n$，从而使得注意力机制天然具备捕捉相对位置关系的能力。
RoPE 被广泛应用于 LLaMA、Qwen 等大语言模型中。

然而，RoPE 的位置编码存在一个固有限制：它只能处理训练时见过的上下文长度。
从公式可以看出，第 $i$ 组维度的旋转角度为 $m\theta_i$，不同频率的维度旋转速度差异极大。
高频维度（$i$ 较大）旋转快，低频维度（$i$ 较小）旋转慢，在训练时的上下文长度内，低频维度可能只旋转了极小的角度，
模型几乎没有见过更大的位置值。当推理时遇到超出训练范围的更长序列时，
低频维度会进入从未探索过的区域，导致模型性能显著下降。

YaRN（Yet another RoPE extensioN）是一种用于扩展 RoPE 上下文长度的方法。
设缩放因子 $s = L' / L$，其中 $L$ 为训练时的上下文长度，$L'$ 为目标长度。
定义第 $i$ 个维度的波长为 $\lambda_i = 2\pi / \theta_i$，并设定两个波长阈值 $\lambda_{\text{low}}$ 和 $\lambda_{\text{high}}$。
YaRN 为每个维度 $i$ 定义一个缩放因子 $s_i$，将旋转角度从 $m\theta_i$ 修改为 $m\theta_i / s_i$：
$$
	s_i =
	\begin{cases}
		1,                                                                                                        & \lambda_i \leq \lambda_{\text{low}}                      \\[6pt]
		s,                                                                                                        & \lambda_i \geq \lambda_{\text{high}}                     \\[6pt]
		1 + (s - 1) \cdot \dfrac{\lambda_i - \lambda_{\text{low}}}{\lambda_{\text{high}} - \lambda_{\text{low}}}, & \lambda_{\text{low}} < \lambda_i < \lambda_{\text{high}}
	\end{cases}
$$

\section{编码器与解码器}
Transformer 的编码器由一组 N 个编码器串联而成，每一个编码器的构造都是相同的，
由自注意力层和前馈网络层组成，每层的输出都经过残差连接和归一化层（Add \& Norm）。
前馈网络由多个全连接层组成，一般采用 ReLU 激活函数。

相较编码器，解码器的注意力层不太一样，它需要关注编码器的输出，因此解码器的注意力层分为两部分：带掩码的注意力层和交叉注意力层。
交叉注意力层都有两个输入：一个来自带掩码的注意力层，另一个是编码器的输出。让我们用 $\boldsymbol{R}$ 来表示编码器的输出值，
用 $\boldsymbol{M}$ 来表示由带掩码的注意力层输出的注意力矩阵，注意力的计算使用矩阵 $\boldsymbol{M}$ 创建查询矩阵 $\boldsymbol{Q}$，
使用 $\boldsymbol{R}$ 创建键矩阵和值矩阵。
在解码器中的带掩码的注意层中，我们对注意力计算中的矩阵进行掩码处理，确保解码器在生成每个单词时只能关注之前的单词。

\begin{figure}[ht]
	\centering
	\includegraphics[width=0.7\textwidth, keepaspectratio]{6-5.pdf}
	\caption{Transformer总结构}
\end{figure}

$$
	\frac{\boldsymbol{Q}_i \cdot \boldsymbol{K}_i^\mathrm{T}}{\sqrt{d_k}} =
	\begin{bmatrix}
		             & \text{<sos>} & \text{我} & \text{很} & \text{好} \\
		\text{<sos>} & 9.125        & -\infty  & -\infty  & -\infty  \\
		\text{我}     & 5.0          & 12.37    & -\infty  & -\infty  \\
		\text{很}     & 7.25         & 5.0      & 10.37    & -\infty  \\
		\text{好}     & 1.5          & 1.37     & 1.87     & 10.0     \\
	\end{bmatrix}
$$

获得解码器的输出之后，我们将其送入线性层和 softmax 层，线性层将生成一个 logit 向量，其大小等于原句中的词汇量，
然后，将 softmax 函数应用于 logit 向量，从而得到概率分布，最后对其进行采样得到最终的文本输出。

\begin{figure}[ht]
	\centering
	% 第一行
	\begin{subfigure}[b]{0.45\textwidth}
		\centering
		\includegraphics[width=\textwidth]{6-4a.pdf}
		\caption{t = 1}
	\end{subfigure}
	\hfill
	\begin{subfigure}[b]{0.45\textwidth}
		\centering
		\includegraphics[width=\textwidth]{6-4b.pdf}
		\caption{t = 2}
	\end{subfigure}

	\vspace{0.5cm} % 行间距

	% 第二行
	\begin{subfigure}[b]{0.45\textwidth}
		\centering
		\includegraphics[width=\textwidth]{6-4c.pdf}
		\caption{t = 3}
	\end{subfigure}
	\hfill
	\begin{subfigure}[b]{0.45\textwidth}
		\centering
		\includegraphics[width=\textwidth]{6-4d.pdf}
		\caption{t = 4}
	\end{subfigure}

	\caption{解码器流程说明}
\end{figure}
